\subsection{Sistemas Distribuídos e Descentralização}

Sistemas centralizados apresentam riscos elevados em ambientes críticos devido à existência de pontos únicos de falha (\textit{Single Points of Failure} - SPOF) \cite{tanenbaum2021redes}. A descentralização distribui a autoridade e o processamento entre múltiplos nós, aumentando a disponibilidade e a resiliência. No HormuzNet, a remoção de bases estáticas e a adoção de uma malha (\textit{mesh}) de \textit{brokers} permite que o sistema continue operacional mesmo que parte da infraestrutura seja comprometida.

\subsection{Ordenação de Eventos e Relógio de Lamport}

Em um sistema descentralizado, não existe um relógio físico global confiável para ordenar eventos ocorridos em diferentes máquinas. Leslie Lamport propôs o conceito de Relógios Lógicos \cite{lamport1978time} para estabelecer uma relação de "aconteceu antes" (\textit{happened-before}). Cada nó mantém um contador local que é incrementado a cada evento e sincronizado através de mensagens. No HormuzNet, o Relógio de Lamport é utilizado para garantir que todos os \textit{brokers} processem as ocorrências de sensores na mesma ordem lógica, permitindo uma tomada de decisão consistente sobre qual drone deve ser despachado.

\subsection{Comunicação em Grupo com UDP Multicast}

O protocolo UDP Multicast permite o envio de datagramas para um grupo de destinatários interessados de forma eficiente \cite{rfc1112}. Diferente do \textit{Unicast}, onde o emissor precisa enviar uma cópia para cada receptor, o \textit{Multicast} utiliza a infraestrutura de rede para replicar o pacote apenas quando necessário. Esta característica é ideal para o HormuzNet, pois permite que sensores publiquem dados uma única vez e todos os \textit{brokers} da malha recebam a informação simultaneamente, garantindo redundância imediata sem sobrecarregar os sensores.

\subsection{Resiliência e Fallback em Conexões TCP}

O Transmission Control Protocol (TCP) fornece um canal bidirecional confiável e orientado à conexão \cite{rfc793}. Para a coordenação de drones, a confiabilidade é essencial para garantir a entrega de ordens de missão. Entretanto, a persistência da conexão TCP em sistemas dinâmicos exige mecanismos de \textit{fallback}. No projeto, drones implementam uma lógica de reconexão automática: ao perderem o contato com seu \textit{broker} primário, tentam se associar a outros nós da malha, garantindo que a frota permaneça sob comando enquanto houver ao menos um \textit{broker} ativo.

\subsection{Concorrência e Paralelismo em Go}

A linguagem Go foi projetada para sistemas concorrentes, utilizando o modelo de \textit{Communicating Sequential Processes} (CSP) \cite{donovan2015go}. Através de \textit{goroutines} (fios de execução leves) e \textit{channels} (canais de comunicação), é possível gerenciar milhares de conexões e eventos simultâneos com baixo custo computacional. O HormuzNet explora esses recursos para processar datagramas UDP, conexões TCP de drones e sessões WebSocket de monitoramento de forma independente e segura, utilizando primitivas de sincronização (\textit{mutexes}) para proteger o estado global compartilhado entre as \textit{goroutines}.
